\documentclass[main.tex]{subfiles}


\begin{document}

%from pagenumber to up
% \phantomsection 
% \hypertarget{toc}{}

 

 

% номер секции вручную
\setcounter{section}{9
}
\addtocounter{section}{-1} 

\section{Сложение световых волн}


Свет~--- это электромагнитные \emph{волны}, поэтому распространение света можно представлять себе как поток волн. (Сказанное далее относится к волнам любой природы, речь о световых волнах идет для определенности.)
%
\begin{law}
\textbf{Принцип суперпозиции волн.} При наложении волн света в некоторой точке пространства их напряженности (индукции) складываются.
\end{law}
%
Имеются два важных случая сложения волн (рис.\,\ref{fig:p190}).

%     \begin{figure}[H]
%     \centering
% \begin{asy}
% settings.outformat="png";
% // settings.prc=true;
% settings.render=10;

% import three;
% import texcolors;
% import x11colors;
% import solids;

% size(6cm,0);
% currentprojection =
% orthographic((5,5,2));


% //styles
% ///////////////////////////////////////////
% DefaultHead3.size=new real(pen p=currentpen) {return 3mm;};
% defaultpen(fontsize(11pt));
% defaultpen(linewidth(0.4pt));
% pen thick=black+2*linewidth(currentpen);
% /////////////////////////////////////////


% //axes
% //draw(O--2.5X,L=Label("$x$",EndPoint,WSW));
% //draw(O--2.3Y,L=Label("$y$",EndPoint,ESE));


% // sin
% real k=1;
% real r=.5;
% real sc=.5;

% triple li(real t){return (0,t,r*sin(k*t));}
% real tmin=-6pi,tmax=0;
% path3 g=graph(li,tmin,tmax,operator..);
% draw(yscale3(sc)*g,red);

% triple li(real t){return (t,0,-r*sin(k*t));}
% real tmin=0,tmax=8pi;
% path3 g=graph(li,tmin,tmax,operator..);
% draw(xscale3(sc)*g,red);


% //dot
% dot((0,-6pi*sc,0),L=Label("$S_1$",black),W,red);
% dot((8pi*sc,0,0),L=Label("$S_2$",black),W,red);
% dot((0,0,0),L=Label("$M$",black),E);


% //dash
% draw((0,0,0)--(0,-6pi*sc,0),dashed+gray);
% draw((0,0,0)--(8pi*sc,0,0),dashed+gray);




% \end{asy}
% \caption{Наложения двух волн в точке пространства}
% \label{fig:p190}
% \end{figure}

\def\sub{.417}
\begin{figure}[H]%
\hfill
\adjustbox{valign=c}{\begin{subfigure}{\sub\linewidth}
\centering
\includegraphics[]{op/21lprinc_pic1.pdf}%
% \caption{При $s_1=s_2$ вторая машина затратила меньше времени.}
% \label{fig:p42sub1}
\end{subfigure}}%
\hfill
\adjustbox{valign=c}{\begin{subfigure}{\sub\linewidth}
\centering
\includegraphics[]{op/21lprinc_pic2.pdf}%
% \caption{При $t_1=t_2$ обе машины прошли одинаковое расстояние.}
% \label{fig:p42sub2}
\end{subfigure}}%
\hfill\mbox{}
\caption{Наложения двух волн в точке пространства}
\label{fig:p190}
\end{figure}


Пусть точечные источники~$S_1$ и~$S_2$ являются \emph{когерентными} (то есть они излучают свет \emph{одинаковой частоты}\footnote{Свет одной частоты (одного цвета) называется \emph{монохроматическим} светом. Белый свет не является монохроматическим: он представляет собой смесь волн с разными частотами (эти частоты отвечают цветам от красного до фиолетового).}  и имеют \emph{постоянную разность фаз}; в данном случае разность фаз равна нулю~--- источники <<колеблются>> одинаково). Волны от этих источников накладываются друг на друга в точке~$M$.

На рис.\,\ref{fig:p190} (слева) разность расстояний $\Delta = S_2M-S_1M$ (\emph{разность хода}) равна длине волны ($\Delta=\lambda$). В этом случае волны в точке~$M$ складываются \emph{в фазе}~--- в точке~$M$ достигается \emph{максимально возможная освещенность} от данных источников.
%
\begin{law}
    \textbf{Условие максимума.} При наложении когерентных волн освещенность в данной точке максимальна, если разность хода равна \emph{целому} числу длин волн:
    \begin{equation}
        \Delta=m\lambda\quad(m=0,1,2,\ldots).
    \end{equation}
\end{law}
%
На рис.\,\ref{fig:p190} (справа) разность хода $\Delta = S_2M-S_1M$ составляет половину длины волны ($\Delta=\lambda/2$). Теперь волны в точке~$M$ складываются \emph{в противофазе}~--- в точке~$M$ достигается \emph{минимально возможная освещенность} от данных источников.
%
\begin{law}
    \textbf{Условие минимума.} При наложении когерентных волн освещенность в данной точке минимальна, если разность хода равна \emph{полуцелому} числу длин волн:
    \begin{equation}
        \Delta=m\lambda+\frac{\lambda}{2}\quad(m=0,1,2,\ldots).
    \end{equation}
\end{law}
%
% Также можно сказать, что минимуму соответствует разность хода <<целое число длин волн плюс половина длины волны>>.
% В той области пространства, где происходит наложение волн двух когерентных источников, наблюдается \textbf{интерференционная картина}~--- фиксированное, не зависящее от времени распределение мест максимумов и минимумов освещенности. В таком случае говорят, что волны \emph{интерферируют} друг с другом, или происходит \emph{интерференция} волн.
%
Если разность хода принимает какое-то иное значение, не равное целому или полуцелому числу длин волн, тогда волны, приходящие в данную точку, создают в ней колебания с некоторой промежуточной освещенностью (больше минимально возможной освещенности, но меньше максимально возможной).

 
\end{document}